\subsection{Discovery, admission, and retention through bank changes}
\label{app:theory-discovery-admission}

Discovery and retention require two separate budgets: enough opportunities
to admit a missing key, and enough persistent loss weight to protect it
after admission. Bank growth connects the two because admitting a key can
also change the objective used to protect earlier keys.

Retention begins only after a verified key enters active memory. Archive
exploration separates discovery from reliable return to a discovered state
\citep{ecoffet2021return}; replay with behavioral cloning supplies a related
empirical mechanism for reducing forgetting \citep{rolnick2019replay}.
Neither result is a theorem that every execution key enters a finite bank
or maintains a positive probability thereafter. We state the additional
admission and energy conditions explicitly.

\paragraph{Actual insertion is the relevant event.}
Fix a prompt and count opportunities according to a specified schedule.
An opportunity can be a prompt-local query or a global update; the resulting
bounds concern that opportunity count. The admission probability includes
verification, capacity, scheduling, and candidate selection. For example,
$G$ conditionally independent samples with unconditional accepted-key probability $p_b$
contain key $b$ with probability $1-(1-p_b)^G$. This is its admission
probability only if the key remains eligible and insertion is guaranteed
whenever it appears. The following argument permits adaptive histories;
its probability foundation is conditional Borel--Cantelli
\citep[Theorem~4.3.4]{durrett2019probability}.

\begin{lemma}[Conditional admission hazards]
\label{ext:lem-admission-hazards}
Fix a key $b$ not initially banked. Let $(\mathcal F_n)_{n\ge0}$ contain the
complete history after each opportunity, and let $\tau_b$ be the first
opportunity at which a verified representative of $b$ is inserted into active
replay memory, with $\tau_b=\infty$ if this never happens. Define
\[
 h_{b,n}=\Pr(\tau_b=n\mid\mathcal F_{n-1}),\qquad
 H_{b,N}=\sum_{n=1}^N h_{b,n}.
\]
The stopped hazard $h_{b,n}$ is zero after an earlier admission. For $a\ge0$,
\begin{equation}
 \Pr(\tau_b>N,\ H_{b,N}\ge a)\le e^{-a},\qquad
 \Pr(\tau_b=\infty,\ H_{b,\infty}=\infty)=0.
 \label{ext:eq-admission-hazards}
\end{equation}
If deterministic $\underline h_{b,n}\in[0,1]$ satisfy
$h_{b,n}\ge\underline h_{b,n}$ on $\{\tau_b\ge n\}$, then
\begin{equation}
 \Pr(\tau_b>N)\le\prod_{n=1}^N(1-\underline h_{b,n})
 \le\exp\!\left(-\sum_{n=1}^N\underline h_{b,n}\right).
 \label{ext:eq-admission-floors}
\end{equation}
For finitely many targets, summing over initially missing keys bounds
non-admission by $N$.
\end{lemma}
\begin{proof}
Set $Z_N=\mathbf1\{\tau_b>N\}\exp(H_{b,N})$ and $Z_0=1$. Earlier admission
is known at time $n-1$, and both $h_{b,n}$ and $Z_{n-1}$ vanish there.
On survival through $n-1$,
\[
 \mathbb E[Z_n\mid\mathcal F_{n-1}]
 =Z_{n-1}e^{h_{b,n}}(1-h_{b,n})\le Z_{n-1}.
\]
Thus $Z_n$ is a nonnegative supermartingale. Markov's inequality proves the
finite-hazard bound. Since $\mathbb E Z_N\le1$, Fatou's lemma excludes a
positive-probability event on which survival continues forever and
$H_{b,N}\to\infty$, which would force $Z_N\to\infty$. With deterministic
hazard floors, conditioning gives
\[
 \Pr(\tau_b>n)\le(1-\underline h_{b,n})\Pr(\tau_b>n-1).
\]
Iteration and a union bound prove the claims without independence across
opportunities.
\end{proof}

\paragraph{Example: turn an admission floor into a budget.}
Suppose a prompt has four possible correct keys and at least four bank slots.
At each opportunity, every missing key is generated with probability at least
$0.04$ and inserted, conditional on appearing, with probability at least
$1/2$. Its actual admission hazard is therefore at least $0.02$. The lemma
gives $4e^{-0.02N}$ as a failure bound, so
$N=\lceil\log(4/0.05)/0.02\rceil=220$ opportunities suffice for at least
$95\%$ probability of full admission. With only two non-evicting slots,
full admission of those four keys is impossible; after the bank fills, the
missing keys' admission hazards become zero despite continued generation.

Positive probability at every opportunity is insufficient. Independent
opportunities with admission probabilities $(n+1)^{-2}>0$ give
\[
 \Pr(\tau=\infty)=\prod_{n=1}^\infty
 \left(1-\frac1{(n+1)^2}\right)=\frac12.
\]
A non-evicting capacity-$K_{\rm cap}$ bank also gives every unbanked key zero admission
hazard once full. For a fixed sampling law with no such gates, discovery
reduces to coupon collection, with incorrect responses acting as null
coupons. \citet[Theorems~2--3]{anceaume2015coupon} show that, at fixed total
non-null mass, uniform independent sampling minimizes collection time in
stochastic order. This comparison concerns a fixed independent sampling
law; the adaptive admission bounds above supply the appropriate statement
when the training history changes sampling or eligibility.

\paragraph{Charge changes in the replay objective.}
A key's admission must preserve a representative with positive loss weight.
The energy can change when bank membership or weights change, so descent
between admissions alone is insufficient. We account for those changes
explicitly, following the general switched-energy viewpoint of
\citet{branicky1998multiple}. The following elementary bound is proved
directly and does not invoke an equilibrium-stability theorem for neural
training.

\begin{lemma}[Retention with a finite switching-energy budget]
\label{ext:lem-switching-retention}
Let $t_0<t_1<\cdots$ be switching times, with finitely many switches on every
bounded time interval. On interval $r$, define
\[
 \begin{aligned}
 F_r(\theta)&=\sum_{j:d_{rj}>0}d_{rj}\ell_j(\theta)-J(\theta),\\
 \ell_j(\theta)&=-\log\pi_\theta(e_j\mid x_j),\qquad
 J(\theta)\le J_{\max}<\infty.
 \end{aligned}
\]
Each interval uses a finite collection of complete verified exemplars with
nonnegative coefficients; zero coefficients denote absent exemplars.
Assume finite initial energy, continuity of parameters at switches,
nonincrease of $F_r$ within interval $r$, and the finite positive-jump budget
\begin{equation}
 A:=\sum_{r\ge1}
   [F_r(\theta(t_r))-F_{r-1}(\theta(t_r))]_+<\infty.
 \label{ext:eq-switching-budget}
\end{equation}
If exemplar $j$ is protected from interval $r_j$ onward and
$d_{rj}\ge\underline d_j>0$ for all $r\ge r_j$, let
$D=F_0(\theta(t_0))+A+J_{\max}$. Then
\begin{equation}
 p_{\theta(t)}(b_j\mid x_j)\ge\pi_{\theta(t)}(e_j\mid x_j)
 \ge\exp(-D/\underline d_j)>0\qquad(t\ge t_{r_j}).
 \label{ext:eq-switching-retention}
\end{equation}
Each likelihood uses the complete verified-response probability under the
sampling law being bounded. For mean-token loss, $d_{rj}=\alpha_{rj}/L_j$,
with fixed length $L_j>0$.
\end{lemma}
\begin{proof}
Induction over interval descents and switch jumps gives
$F_r(\theta(t))\le F_0(\theta(t_0))+A$. Consequently
$\sum_jd_{rj}\ell_j(\theta(t))\le D$. Each summand is nonnegative, so a
protected exemplar satisfies $\ell_j(\theta(t))\le D/\underline d_j$.
The key event contains its exemplar event, so exponentiation proves
the bound.
\end{proof}

If admission preserves old coefficients, the jump equals the weighted
surprisal of the newly added exemplars at their admission probabilities.
Uniform categorical replay with fixed total dose $\rho$ instead dilutes old
weights. For a size-$k$ bank, write
$R_k=-k^{-1}\sum_{b\in\mathcal B_k}\log p_b$; adding one new key gives
\[
 \Delta F=\frac{\rho}{k+1}[-\log p_{\rm new}-R_k].
\]
Old coefficients decrease from $\rho/k$ to $\rho/(k+1)$ but stay at least
$\rho/K_{\rm cap}$ under capacity $K_{\rm cap}$. Finitely many admissions with positive
admission-time probabilities and no eviction therefore have finite total
switching cost in the exact descent model. Changing a representative must
also be charged to the objective jump. A key-level version of
Lemma~\ref{ext:lem-switching-retention} follows by requiring at least one
complete representative of that key with coefficient bounded below on every
subsequent interval; the same nonnegative-summand proof applies even if the
representative changes.

\paragraph{Example: adding a rare key has a calculable energy cost.}
Take $\rho=0.3$ and a two-key categorical bank whose keys each have
probability $1/4$, so $R_2=\log4$. Add a third key with probability $1/16$,
without changing the policy at that instant. Then
\[
 \Delta F=\frac{0.3}{3}(\log16-\log4)
          =0.1\log4\approx0.1386.
\]
The earlier keys' coefficients decrease from $0.15$ to $0.10$. A
capacity-three bank still supplies a $0.10$ coefficient floor, but the
positive jump $0.1386$ must enter the energy budget. This is why retention
through admission needs both the coefficient floor and the switch cost.

\begin{corollary}[Conditional discovery followed by retention]
\label{ext:cor-discovery-retention}
Suppose one prompt's entire correct support $\mathcal C$ is finite, of size
$m$ no larger than bank capacity. Admissions do not evict keys, and every
missing key $b$ has conditional admission hazard at least $\epsilon r_b$ at
every opportunity, for constants $\epsilon>0$, $r_b>0$, and
$\epsilon r_b\le1$. Assume the switching-energy conditions of
Lemma~\ref{ext:lem-switching-retention}, including a finite positive-jump
budget and persistent positive coefficients for admitted complete exemplars.
Then full correct coverage is reached in finitely many opportunities almost
surely. Thereafter each correct key has a uniform-in-time positive
probability floor. At horizon $N$,
\begin{equation}
 \Pr(\text{full coverage by }N)
 \ge 1-\sum_{b\in\mathcal C}e^{-N\epsilon r_b}.
 \label{ext:eq-discovery-retention}
\end{equation}
If every $r_b\ge\gamma>0$, this is at least $1-me^{-N\epsilon\gamma}$.
\end{corollary}
\begin{proof}
Apply Lemma~\ref{ext:lem-admission-hazards} to the initially missing keys and
take a union bound. The failure bound tends to zero. Finite support and
non-eviction imply a finite almost-sure opportunity index of simultaneous
full coverage. Lemma~\ref{ext:lem-switching-retention} protects each admitted
complete exemplar and hence its key; after full coverage this applies to
all correct keys.
\end{proof}

A fixed full-support proposal mixture supplies the assumed hazard floor if
it is queried at every counted opportunity and insertion is guaranteed
whenever a missing key appears. Capacity must accommodate the entire
admissible correct support, including keys competing for the same slots.
Finite-temperature sampling, bounded search attempts, or recurrent replay
scheduling alone do not establish this admission floor or the energy
budget. The probability bound~\eqref{ext:eq-discovery-retention} controls
admission by a finite opportunity count; the subsequent retention floors
can depend on the realized training history. The corollary establishes
neither uniform key probabilities nor convergence of the neural optimizer.

